\documentclass{article}
\usepackage[T1]{fontenc}
\usepackage{lmodern}
\usepackage{graphicx}
\usepackage{amsmath,amssymb}
\usepackage[a4paper,margin=2cm]{geometry}
\usepackage[absolute,overlay]{textpos}
\setlength{\TPHorizModule}{1mm}
\setlength{\TPVertModule}{1mm}
\usepackage[indonesian]{babel}
\usepackage{hyperref}
\setlength{\emergencystretch}{2em}
% unit: Halaman 18

\begin{document}

% === Halaman 18 (llm) ===

\begin{itemize}
    \item Jika panjang plainteks tidak habis dibagi dengan ukuran blok, maka blok terakhir berukuran lebih pendek daripada blok-blok lainnya.
    \item Untuk itu, kita tambahkan bit-bit \textit{padding} untuk menutupi kekurangan bit blok.
    \item Bit-bit padding misalnya bit 0 semua, atau bit 1 semua, atau bit 0 dan bit 1 berselang-seling.
    \item Contoh: Ukuran blok 8 bit: \textcolor{red}{10001101} \textcolor{red}{00101001} \textcolor{red}{10110}\textcolor{blue}{000}\\
    Blok terakhir hanya 5 bit, ditambah 3 bit \textit{padding} (\textcolor{blue}{warna biru})
\end{itemize}

\end{document}
